\section{Mokamint's Proof of Space Mining Algorithm}\label{sec:pos}

Mokamint uses a mining algorithm based on a specific instance of proof of space,
derived from Signum's. The idea is relatively simple: each peer
allocates, on disk, its own \emph{plot}, \ie{}, a set of data structures called \emph{nonces}.
A nonce (Fig.~\ref{fig:waiting_time_nonce}) is an array of bytes, split into chunks called \emph{scoops}.
Whenever a peer needs to create a new block, to expand a chain of blocks $\xi$,
it computes a \emph{challenge} from $\xi$ that sends to its miner (Fig.~\ref{fig:mining}). The miner uses
that challenge to compute the minimal waiting time among that of the nonces
in its plot on disk (Fig.~\ref{fig:waiting_time_plot}).
The nonce that generates the minimal
waiting time is sent to the peer, that will selevct the minimal waiting time among those
of all its connected miners and use it to build the block that expands $\xi$.
That block can then be propagated to all peers, but only after the waiting time has elapsed,
or otherwise other peers might reject it as (still) invalid.

\begin{figure}[t]
  \begin{center}
    \begin{tikzpicture}[scale=1]

      \draw (-5.5,1) node {nonce};
      \draw [thick] (-11.5,0) rectangle +(1.75,0.5);
      \draw (-10.6,0.2) node {$\scoop_0$};
      \draw [thick] (-9.75,0) rectangle +(1.75,0.5);
      \draw (-8.85,0.2) node {$\scoop_1$};
      \draw [thick] (-8,0) rectangle +(1.75,0.5);
      \draw (-7.1,0.2) node {$\scoop_2$};
      \draw [dashed] (-6.25,0) rectangle +(1,0.5);
      \draw (-5.75,0.25) node {$\ldots$};
      \draw [thick,fill=verylightred] (-5.25,0) rectangle +(1.75,0.5);
      \draw (-4.375,0.2) node {$\scoop_i$};
      \draw [dashed] (-3.5,0) rectangle +(1,0.5);
      \draw (-3,0.25) node {$\ldots$};
      \draw [thick] (-2.5,0) rectangle +(1.75,0.5);
      \draw (-1.63,0.2) node {$\scoop_{4094}$};
      \draw [thick] (-0.75,0) rectangle +(1.75,0.5);
      \draw (0.13,0.2) node {$\scoop_{4095}$};

      \draw (-5.5,-0.5) node {the waiting time of the nonce \wrt{} a challenge $\langle i,\gamma\rangle$ is $\hash(\scoop_i\append\gamma)$};
%      \draw (-5.5, -0.9) node {normalized \wrt{} the current acceleration of the network};

    \end{tikzpicture}
  \end{center}
  \caption{A nonce is an array of bytes, split into chunks called scoops.
    The computation of the waiting time of a nonce \wrt{} the challenge
    $\langle i,\gamma\rangle$ selects the $i$th scoop of the nonce
    and computes a hash value that is interpreted as a numerical waiting time.}\label{fig:waiting_time_nonce}
\end{figure}

Let us formalize these mining algorithm now.
%
\begin{definition}[Challenge]\label{def:challenge}
  The set of \emph{challenges} is
  \[
  \Challenges=\left\{\langle\scoopnumber,\gamma\rangle\left|
  \begin{array}{l}
    0\le\scoopnumber<4096\text{ and $\gamma$ is a hash}
  \end{array}
  \right.\right\}.
  \]
  The $\gamma$ component of a challenge is said to be its \emph{generation}.
\end{definition}
%
The intuition is that $\gamma$ is a sort of digest of the history of the chain $\xi$ that
one wants to expand with a next block. Sec.~\ref{sec:blockchain_construction} will
show how a challenge is derived from $\xi$.

A \emph{scoop} is a pair of hashes. A \emph{nonce} is a \emph{progressive number} $p$, and
a list of $4096$ scoops (see Fig.~\ref{fig:waiting_time_nonce}, where the progressive
number is not shown, for simplicity).
%
\begin{definition}[Scoop, Nonce]\label{def:scoop_nonce}
  The sets of \emph{scoops} and \emph{nonces} are
  \begin{align*}
    \Scoops&=\left\{\langle h_1,h_2\rangle\left|\begin{array}{l}
    h_1,h_2\text{ are hashes}
    \end{array}\right.\right\},\\
    \Nonces&=\left\{\langle p,\scoops\rangle\left|\begin{array}{l}
    p\in\nat\text{ and }\scoops\in\Scoops^{4096}
    \end{array}\right.\right\}.
  \end{align*}
\end{definition}

\begin{figure}[t]
  \begin{center}
    \begin{tikzpicture}[scale=1]

      \draw (0.4,5.4) node {plot};

      \draw [thick] (-0.5,4.5) rectangle +(1.75,0.5);
      \draw (0.4,4.75) node {$\nonce_0$};
      \draw[->] (1.5,4.75) -- +(0.25,0);
      \draw (4.3,4.75) node {waiting time of $\nonce_0$ \wrt{} $\langle i,\gamma\rangle$};

      \draw [thick] (-0.5,4) rectangle +(1.75,0.5);
      \draw (0.4,4.25) node {$\nonce_1$};
      \draw[->] (1.5,4.25) -- +(0.25,0);
      \draw (4.3,4.25) node {waiting time of $\nonce_1$ \wrt{} $\langle i,\gamma\rangle$};

      \draw [thick,fill=verylightred] (-0.5,3.5) rectangle +(1.75,0.5);
      \draw (0.4,3.75) node {$\nonce_2$};
      \draw[->] (1.5,3.75) -- +(0.25,0);
      \draw (4.3,3.75) node {waiting time of $\nonce_2$ \wrt{} $\langle i,\gamma\rangle$};

      \draw [thick] (-0.5,3) rectangle +(1.75,0.5);
      \draw (0.4,3.25) node {$\nonce_3$};
      \draw[->] (1.5,3.25) -- +(0.25,0);
      \draw (4.3,3.25) node {waiting time of $\nonce_3$ \wrt{} $\langle i,\gamma\rangle$};

      \draw (0.4,2.375) node {$\vdots$};
      \draw [dashed] (-0.5,1.5) rectangle +(1.75,1.5);
      \draw (4.3,2.375) node {$\vdots$};

      \draw [thick] (-0.5,1) rectangle +(1.75,0.5);
      \draw (0.4,1.25) node {$\nonce_{n-3}$};
      \draw[->] (1.5,1.25) -- +(0.25,0);
      \draw (4.475,1.25) node {waiting time of $\nonce_{n-3}$ \wrt{} $\langle i,\gamma\rangle$};

      \draw [thick] (-0.5,0.5) rectangle +(1.75,0.5);
      \draw (0.4,0.75) node {$\nonce_{n-2}$};
      \draw[->] (1.5,0.75) -- +(0.25,0);
      \draw (4.475,0.75) node {waiting time of $\nonce_{n-2}$ \wrt{} $\langle i,\gamma\rangle$};

      \draw [thick] (-0.5,0) rectangle +(1.75,0.5);
      \draw (0.4,0.25) node {$\nonce_{n-1}$};
      \draw[->] (1.5,0.25) -- +(0.25,0);
      \draw (4.475,0.25) node {waiting time of $\nonce_{n-1}$ \wrt{} $\langle i,\gamma\rangle$};

      \draw [decorate,decoration={brace,amplitude=4pt},thick,xshift=0.1cm,yshift=0pt]
      (7.2,5) -- (7.2,0) node [midway,right,xshift=.2cm] {$\min$};

    \end{tikzpicture}
  \end{center}
  \caption{A plot is a set of nonces. Its waiting time
    \wrt{} a challenge $\langle i,\gamma\rangle$ is the minimal waiting time among those of its nonces
    \wrt{} $\langle i,\gamma\rangle$ (Fig.~\ref{fig:waiting_time_nonce}).
    For instance, it might be that of $\nonce_2$ in the picture.}\label{fig:waiting_time_plot}
\end{figure}

A \emph{prolog} identifies the parts involved in mining: the peer, the miner (Fig.~\ref{fig:genericity})
and the chain identifier of the blockchain for which mining is performed.
%
\begin{definition}[Prolog]\label{def:prolog}
  The set of \emph{prologs} is
  \[
  \Prologs=\left\{\langle\mathit{key}_\public^{\mathit{peer}},
  \mathit{key}_\public^{\mathit{miner}},\chainid\rangle\left|\begin{array}{l}
  \text{where $\mathit{key}_\public^{\mathit{peer}}$
    and $\mathit{key}_\public^{\mathit{miner}}$ are the public keys}\\
  \text{that identify a peer and a miner, respectively,}\\
  \text{and $\chainid$ is the chain identifier of a blockchain}
  \end{array}\right.\right\}.
  \]
\end{definition}
%
Prologs are relevant since they are used to mark nonces:
each nonce and consequently each plot is \emph{for a given prolog}.
Therefore, a miner cannot use the same plot for mining for more peers or for more
chains.

Def.~\ref{def:nonce_construction} formalizes the construction of a nonce, given its progressive number and a prolog. It uses a constant $\kappa>0$
to limit its computational cost, by avoiding hashing very large chunks of data.
The construction algorithm derives a sequence of hashes (steps~\ref{step:nonce_construction:seed}
and~\ref{step:nonce_construction:first_hash}) and constructs a \emph{final} hash from
all of them (step~\ref{step:nonce_construction:final_hash}) that uses to modify
the original sequence of hashes (step~\ref{step:nonce_construction:modify}).
Then it shuffles the sequence (step~\ref{step:nonce_construction:swap}), which intuitively
guarantees that, in order to compute a given scoop of the nonce
(\ie{}, a pair of consecutive hashes), the algorithm must be thoroughly executed.
%
\begin{definition}[$\nonce(p,\pi)$]\label{def:nonce_construction}
  Given $p\in\mathbb{N}$ and $\pi\in\Prologs$, let
  \[
  \nonce(p,\pi)=\langle p,\langle h_0,h_1\rangle\append
  \cdots\append\langle h_{2\cdot 4096-2},h_{2\cdot 4096-1}\rangle\rangle\in\Nonces,
  \]
  where the hashes $h_0,\ldots,h_{2\cdot 4096-1}$ are constructed as follows:
  %
  \begin{enumerate}
  \item\label{step:nonce_construction:seed} let $\seed=\pi.\mathit{key}^{\mathit{peer}}_{\public}\append
    \pi.\mathit{key}^{\mathit{miner}}_{\public}\append\pi.\mathit{chainid}\append p$;
  \item\label{step:nonce_construction:first_hash}
    for each $i$ from $2\cdot 4096-1$ to $0$, let
    \[
    h_i=\hash\left(\text{first $\kappa$ bytes of}\left(\left(\append_{i<j<2\cdot 4096}h_j\right)\append\seed\right)\right);
    \]
  \item\label{step:nonce_construction:final_hash}
    let
    $
    \finalhash=\hash\left(\left(\append_{0\le j<2\cdot 4096}h_j\right)\append\seed\right);
    $
  \item\label{step:nonce_construction:modify} for each $i$ from $0$ to $2\cdot 4096-1$, reassign
    $h_i$ to $h_i\xor\finalhash$;
  \item\label{step:nonce_construction:swap}
    for each odd $i$ from $1$ to $2\cdot 4096-1$, swap
    $h_i$ with $h_{2\cdot 4096-i}$.
  \end{enumerate}
\end{definition}
%
Let us discuss Def.~\ref{def:nonce_construction} now.
Step~\ref{step:nonce_construction:seed} starts a chain of hash computations that includes the prolog $\pi$
in the seed. As a consequence, the resulting nonces will be \emph{for $\pi$}: for instance,
the nonces of another peer or of another miner or for another chain identifier
will have a different $\pi$ and consequently
they will be different nonces.
Step~\ref{step:nonce_construction:first_hash} is a loop that computes all $h_i$ from the seed.
It uses a hashing function $\hash$ that is expected to be expensive to compute, such
as shabal256: this makes it hard to recompute nonces and plots on-the-fly (which would
be a proof of work algorithm) rather than once and for all before mining.
Step~\ref{step:nonce_construction:final_hash} computes a digest $\finalhash$ of all $h_i$'s and
step~\ref{step:nonce_construction:modify} uses that digest to make each $h_i$ dependent on all others, so that
one must compute them all in order to compute each of them.
This is strengthened by the traditional use of $\oplus$ (the bitwise xor function) to minimize the information loss.
Step~\ref{step:nonce_construction:swap} makes each scoop include a $h_i$ with $i\ge 4096$,
so that smaller scoops are not easier to compute than the others. This avoids a potential
time/memory tradeoff attack~\cite{ParkKFGAP18}.

Mokamint miners hold one or more plots on disk. Their initialization is performed
only once and offline, hence it is not part of the mining protocol.
Namely, a plot is a set of nonces constructed according to Def.~\ref{def:nonce_construction},
for a finite non-empty set of progressive numbers $P$ and
for a given prolog $\pi$.
%
\begin{definition}[Plot]\label{def:plot}
  The set of \emph{plots} is defined as
  \[
  \Plots=\left\{\left<\pi,\nonces\right>\left|\begin{array}{l}
  \pi\in\Prologs,\ \varnothing\not=P\subset\mathbb{N}\text{ is finite and }
  \nonces=\left\{\nonce(p,\pi)\left|\;p\in P\right.\right\}
  \end{array}\right.\right\}.
  \]
\end{definition}

Given a challenge and a nonce, the latter has a \emph{waiting time} that specifies how well
the nonce answers the challenge (Fig.~\ref{fig:waiting_time_nonce}).
%
\begin{definition}[$\mathit{waitingTime}(\nonce,\challenge)$]\label{def:nonce_value}
  Let $\nonce\in\Nonces$ and $\challenge\in\Challenges$.
  The \emph{waiting time of $\nonce$ \wrt{} $\challenge$} is
  \[
    \mathit{waitingTime}(\nonce,\challenge)=\hash(\nonce.\scoops[\challenge.i]\append\challenge.\gamma).
  \]
\end{definition}
%
Miners use the waiting time of their nonces
to choose the best nonce for each given challenge,
\ie{}, the one with minimal value (Fig.~\ref{fig:waiting_time_plot}).
Note that only the $(\challenge.i)\mathit{th}$ scoop is used to compute that value.
That is, Mokamint consults only a small portion of the whole plot
file at each mining challenge.
That nonce could then be sent to the peer as answer for
the challenge (Fig.~\ref{fig:mining}).
But nonces are relatively big: around $262$ kbytes. Since they
must be stored in blockchain, inside each block, and since they should be
exchanged between miner and peer
(Fig.~\ref{fig:mining}), a more compact data structure is needed.
\emph{Deadlines} are exactly this compact data structure
that represents the nonce computed for a challenge.
Namely, a deadline carries the information needed
to reconstruct the nonce (progressive $p$ and prolog $\pi$, see Def.~\ref{def:nonce_construction}),
the value of that nonce (to avoid its repeated expensive recomputation through the formula
in Def.~\ref{def:nonce_value}), and
the challenge that that nonce was meant to answer. Therefore, in Fig.~\ref{fig:genericity}, the miner
will actually send back deadlines rather than nonces to the peer, which reduces the
data throughput. This is paramount if peer and miner are connected through a network rather
than being in the same physical machine.
%
\begin{definition}[Deadline]\label{def:deadline}
  The set of \emph{deadlines} is
  \[
  \Deadlines=\left\{
  \langle p,\pi,\mathit{value},\challenge\rangle
  \left|\begin{array}{l}
  \pi\in\Prologs,\ p\in\mathbb{N},\ \mathit{value}\text{ is a hash}\\
  \text{and }\challenge\in\Challenges
  \end{array}
  \right.
  \right\}.
  \]
\end{definition}
%
When a miner has selected its best nonce for a challenge, from a plot with prolog $\pi$,
and wants to send it to the peer,
it actually sends the much smaller deadline $\delta(\nonce,\pi,\challenge)$.
%
\begin{definition}[$\delta(\nonce,\pi,\challenge)$]\label{def:deadline_from_nonce}
  Given $\nonce\in\Nonces$, $\pi\in\Prologs$ and $\challenge\in\Challenges$, the
  \emph{deadline computed from $\nonce$ for $\pi$ and $\challenge$} is
  $\delta(\nonce,\pi,\challenge)=\langle\nonce.p,\pi,\mathit{value}(\nonce,\challenge),\challenge\rangle$.
\end{definition}

\begin{algorithm}[t]
  \Input{a $\plot$, already computed on disk}
  \Input{the peer $P$ connected to $M$}
  \BlankLine
  \thread{$\mathsf{mineDeadline}(\plot,P)$}{
    \forever{}{
      \nl\label{step:deadline_mining:wait}
      wait to receive a $\challenge\in\Challenges$ from $P$\;
      %
      \nl\label{step:deadline_mining:minimize}
      identify the $\nonce\in\plot.\nonces$ that minimizes its waiting time
      \wrt{} $\challenge$ (Def.~\ref{def:nonce_value})\;
      %
      \nl\label{step:deadline_mining:compute_deadline}
      send $\delta(\nonce,\plot.\pi,\challenge)$ to $P$ (Def.~\ref{def:deadline_from_nonce})\;
    }
  }
  \caption{The deadline mining algorithm, performed by each miner $M$ in its own thread.}
  \label{alg:deadline_mining}
\end{algorithm}

Alg.~\ref{alg:deadline_mining} is executed, inside a dedicated thread,
by each Mokamint miner $M$, connected to a peer $P$. It is an infinite loop that
waits for a challenge (step~\ref{step:deadline_mining:wait}), finds the best nonce
for that challenge, among those in the precomputed plot stored on its disk
(step~\ref{step:deadline_mining:minimize}), and sends back to $P$ the deadline
corresponding to that nonce (step~\ref{step:deadline_mining:compute_deadline}).

As already observed, step~\ref{step:deadline_mining:minimize}
of Alg.~\ref{alg:deadline_mining}, for each received challenge,
consults only a small portion of the plot file to identify the
nonce with minimal waiting time:
it consults a single scoop for each nonce (Fig.~\ref{fig:waiting_time_nonce}).
Therefore, that step requires a fraction of
second only, even for large plot files. It follows that the life of a miner consists almost only in waiting
at step~\ref{step:deadline_mining:wait}, and that a miner is almost always idle
(no CPU computation, no disk activity), which explains the minimal energy-consumption of
Alg.~\ref{alg:deadline_mining}.
Namely, Fig.~\ref{fig:plot_and_mining_costs} reports the time for the creation of a plot file and the time for
using that file for mining, collected with Mokamint. Both grow linearly with the number of nonces in the plot file.
We recall that the goal of Signum's proof of space is to make the former expensive and the latter cheap. For instance, Fig.~\ref{fig:plot_and_mining_costs} shows that,
even with a powerful CPU, a plot of $4000$ nonces cannot
be recomputed at every block challenge, if the block creation rate is $20$ seconds.
That is, the nonces must be precomputed and stored on disk.

\begin{figure}[t]
  \begin{tabular}{|c||c|c|c|c|c|c|c|}
    \hline
    nonces & size & time for creation & time per challenge \\\hline\hline
    1000 & 256MBytes & 7.85s & 0.00078s \\\hline
    2000 & 512MBytes & 14.73s & 0.00156s \\\hline
    4000 & 1GByte & 31.47s & 0.00273s \\\hline
    8000 & 2GBytes & 64.80s & 0.00546s \\\hline
    16000 & 4GBytes & 130.99s & 0.01131s \\\hline
    32000 & 8GBytes & 257.71s & 0.02223s \\\hline
    64000 & 16GBytes & 554.25s & 0.04407s \\\hline
  \end{tabular}
  \caption{Time for the creation of plot files and for their use for mining:
    \emph{nonces} is the number of nonces in the plot file; \emph{size} is the size
    of the plot file; \emph{time for creation} is the time (in seconds) for creating the
    plot file (Def.~\ref{def:plot} and~\ref{def:nonce_construction});
    \emph{time for challenge} is the time (in seconds) to answer a challenge
    to compute a deadline, with that plot file (steps~\ref{step:deadline_mining:minimize}
    and~\ref{step:deadline_mining:compute_deadline} of Alg.~\ref{alg:deadline_mining}).
    Experiments have been performed on am Intel Core i7-11800H CPU.
    They can be verified by executing the \texttt{PlotProfiler} class
    of the \texttt{io-mokamint-plotter} module of~\cite{MokamintSourceCode}.
  }
  \label{fig:plot_and_mining_costs}
\end{figure}
